\subsection{Exploratory Recognition Experiments}
\label{sec:recognition-experiments}

The baseline results above describe the original implementation. Two subsequent
experiments investigated whether the extraction mask and recognition parameters
could improve these results. Both experiments were implemented in a separate
\code{experiments/} directory and executed using the independent
\code{BallExperiments} target. They did not replace the classification functions,
configuration values, or execution path of the main application. The results in
this subsection therefore describe experimental alternatives rather than a
change to the deployed baseline.

\subsubsection{Evaluation protocol and interpretation}

For these experiments, Table1--Table6 were used as an exploratory tuning set.
AllBalls was reported separately as a reference image and was not included in
the ranking objective of the recognition grid search. This differs from the
earlier detection experiment, which selected Hough parameters on AllBalls and
evaluated them on Table1--Table6. Since the recognition search used the labels
of Table1--Table6 to select thresholds, its final scores on those images are
tuning scores, not independent test estimates. Furthermore, all seven images
share the same viewpoint and some unchanged ball positions. AllBalls is also
the image on which the original rules were developed. Its separate results
are consequently a useful regression check, but not evidence of generalisation
to a new camera, table or lighting condition.

The detector, pocket filtering, perspective transformation and matching
tolerance were held fixed. Detection was performed once per image. Predicted
centres were transformed into table coordinates and matched one-to-one to
annotated ball centres with a maximum distance of 15 table pixels, using the
existing evaluation implementation. A match was established before varying
recognition parameters and reused for every configuration. The experimental
run contained 90 annotated balls in Table1--Table6, of which 86 were matched.
There were two unmatched candidate detections, which were excluded from
recognition accuracy and displayed with an unavailable correctness indicator.
All 16 annotated balls in AllBalls were matched. These are counts measured
by the experiment's actual detection run; they are not inferred from a
previously saved detection table.

For a matched ball, its annotated number determines its true type: number 0
denotes the white cue ball, numbers 1--8 are solid, and numbers 9--15 are
striped. Candidate indices identify detections within an image and should not
be interpreted as ball numbers.

Four measures were recorded. Let $M$ be the number of matched detections,
$G$ the number of annotated balls, $t_i$ the true type and $n_i$ the true
number of matched candidate $i$. Let $\hat{t}_i$ and $\hat{n}_i$ be the
predicted type and number. The recognition accuracies are
\[
 A_{\mathrm{type}}=\frac{1}{M}\sum_{i=1}^{M}
       \mathbf{1}(\hat{t}_i=t_i),\qquad
 A_{\mathrm{number}}=\frac{1}{M}\sum_{i=1}^{M}
       \mathbf{1}(\hat{n}_i=n_i).
\]
The number prediction uses the predicted type, so a type error can propagate
into number recognition. An additional diagnostic reruns number classification
with the true type supplied:
\[
 A_{\mathrm{oracle}}=\frac{1}{M}\sum_{i=1}^{M}
       \mathbf{1}\bigl(f_{\mathrm{number}}(\mathrm{clusters}_i,t_i)=n_i\bigr).
\]
This diagnostic helps separate colour-selection failures from type errors.
It is not an executable recognition solution because ground truth is required.
In particular, supplying the true type of a cue ball invokes the existing
number-0 shortcut when segmentation is available. Thus the measure should not
be described as a pure hue-classification accuracy. Finally, end-to-end number
recall divides correct detected-and-numbered balls by $G$ instead of $M$.
Missed balls remain errors under this measure. It does not measure the
precision of numbered detections, since false-positive candidates are not
included in its numerator or denominator.

\subsubsection{Radius padding and cloth removal}

The initial hypothesis was that Hough detection could underestimate a ball's
radius and exclude useful pixels around its edge. The extraction radius was
therefore changed from the rounded Hough radius $r$ to
$r'=r+p$, where $p\in\{0,1,2,3,4\}$ pixels. Both the square crop and its
circular mask used $r'$. Enlarging the square alone would not admit additional
pixels into colour analysis, because pixels outside the original mask would
still be excluded. Crops were clipped to the image boundaries. Padding zero
without filtering reused the original extraction function and reproduced the
baseline mask.

Each padding value was evaluated with and without a local cloth-removal
filter, producing ten variants per candidate. K was fixed at 2 and all type
and number thresholds retained their original values. This isolated the effect
of extraction from changes to the recognition rules.

The filter estimated the surrounding cloth colour from an annulus of width
six pixels immediately outside the padded circle. Sampling used the original
colour image. A sample was eligible when its OpenCV HSV values satisfied
$35\leq H\leq95$, $S\geq50$ and $V\geq30$. These conditions restricted
the estimate to visible green material rather than including arbitrary
colours from nearby objects. If fewer than 20 eligible samples were available,
filtering was skipped and the padded circular mask was retained.

The modal hue bin of eligible samples was selected. Saturation and brightness
were averaged over samples within ten hue units of that mode. Pixels inside
the candidate's mask were marked as potential cloth when their hue differed
from the modal hue by at most 10, their saturation and brightness differed
from the corresponding averages by at most 70, and they satisfied the same
minimum saturation and brightness conditions. These tolerances use OpenCV's
$0$--$179$ hue and $0$--$255$ saturation/brightness scales.

Potential cloth pixels were divided into eight-connected components. Only
components touching the boundary of the circular mask were removed. The
boundary was computed by subtracting a $3\times3$ elliptical erosion from
the original mask, with pixels outside the crop treated as zero. This also
recognised image-clipped edges as boundaries. Removed pixels were excluded
from the mask rather than replaced by black pixels inside the analysis area.
Consequently, K-means percentages used only the retained pixels as their
denominator. The black background in diagnostic images was for display.

Boundary connectivity reduces the chance of deleting isolated green markings
inside another ball, but does not distinguish a green ball from the cloth
when their colours are similar and the matching region reaches the boundary.
The observed masks showed that real green ball regions could be removed.
Conversely, shadows and reflections could prevent some cloth pixels from
matching the estimated colour. Cloth removal was therefore treated as an
experimental option rather than assumed to be beneficial.

% Generated by experiments/writeRecognitionReport.py from saved CSVs.
\begin{table}[htbp]
\centering
\small
\caption{Radius-padding experiment on matched balls in Table1--6 using the original recognition thresholds and K=2.}
\label{tab:radius-padding-results}
\begin{tabular}{rcrrr}
\toprule
Padding (px) & Cloth filter & Matched & Correct type & Correct number \\
\midrule
0 & Off & 86 & 76 & 67 \\
1 & Off & 86 & 73 & 60 \\
2 & Off & 86 & 67 & 49 \\
3 & Off & 86 & 62 & 43 \\
4 & Off & 86 & 60 & 42 \\
0 & On & 86 & 77 & 69 \\
1 & On & 86 & 76 & 66 \\
2 & On & 86 & 78 & 64 \\
3 & On & 86 & 72 & 55 \\
4 & On & 86 & 66 & 50 \\
\bottomrule
\end{tabular}
\end{table}
\FloatBarrier


Without filtering, increasing padding reduced correct number predictions on
Table1--Table6 from 67 at zero padding to 60, 49, 43 and 42 at one through
four pixels respectively. Cloth filtering partially compensated for the
additional background, but padded variants still failed to improve number
recognition over the original baseline. Two-pixel padding with filtering
increased correct type predictions from 76 to 78, but reduced correct numbers
from 67 to 64. The zero-padding filter alone gave 77 correct types and 69
correct numbers. On AllBalls, the unfiltered baseline already classified all
16 balls correctly; the filter at zero padding reduced both counts to 15.

These findings support retaining the original extraction radius for the next
experiment. A likely explanation for the deterioration with padding is that
additional cloth pixels alter cluster colours and proportions. This is an
interpretation of the masks and recognition changes, not a direct measurement
of the geometric accuracy of the Hough radius. The results do not establish
that every detected radius equals the true visible ball boundary.

\subsubsection{Recognition search: fixed processing and caching}

The next experiment varied K, the interpretation of cluster colours and the
type and number decision thresholds. Radius padding was fixed at zero.
Cloth filtering remained an on/off factor, using the same filter parameters
as the preceding experiment. Connected-component type correction was not
included, so changes could be attributed to clustering and colour rules.

K-means was performed in three-dimensional BGR space using the existing
implementation. The requested cluster counts were $K\in\{2,3,4,5,6\}$.
For each call, the random-number generator state was set to 12345, three
K-means++ initialisations were attempted, and the termination criteria were
50 iterations or an epsilon of 0.2. The prior random state was restored
afterwards. These controls make comparisons repeatable within the same
software environment; bit-for-bit equality across different OpenCV versions
or platforms is not assumed.

Each candidate was segmented once for each K/filter combination. The run
contained 104 candidates across all images, giving 1,040 cached segmentations.
Only cluster percentages and representative BGR/HSV colours were retained
for threshold evaluation. Cluster masks and segmented images were released
after caching because subsequent rules depend on these features rather than
on individual pixel labels. The original masked crops were retained for
final diagnostic figures. No K-means rerun was required for each threshold
configuration.

The experimental classifiers accepted explicit parameter structures rather
than changing constants in the main classifier. Before searching, the
experimental default type, percentages and number predictions were compared
with the original production functions on every cached segmentation. Any
disagreement terminated the experiment. This checked that baseline scores
and parameterised scores represented the same original decision rules.

\subsubsection{Stage one: joint K and type-threshold grid}

The first stage evaluated the full Cartesian product below. The original
number thresholds were fixed during this stage, although changing the light
thresholds also affected which cluster was eligible for number recognition.
The same light definition was used in type classification and body-colour
selection.

\begin{table}[htbp]
\centering
\small
\caption{Stage-one recognition search space.}
\label{tab:recognition-search-space}
\begin{tabular}{ll}
\toprule
Parameter & Evaluated values \\
\midrule
Number of clusters K & 2, 3, 4, 5, 6 \\
Cloth filtering & Off, on \\
Minimum light-cluster V & 150, 170, 190 \\
Maximum light-cluster S & 125, 145, 165, 185 \\
Minimum light proportion for cue ball & 80\%, 85\%, 90\% \\
Minimum non-light proportion for solid ball & 60\%, 65\%, 70\%, 75\%, 80\% \\
\bottomrule
\end{tabular}
\end{table}
\FloatBarrier

This produced $5\times2\times3\times4\times3\times5=1{,}800$
configurations. A cluster was light when $V$ met the minimum and $S$ did
not exceed the maximum. A ball was first considered for the cue-ball rule,
then the solid rule, and finally the striped rule. Percentages and ordering
therefore followed the original implementation. A failed segmentation or
absence of sufficient evidence remained unknown and counted as incorrect
for a matched ball.

The dark-cluster threshold was fixed at $V\leq90$. Although it changes the
separate reported dark and coloured percentages, it does not change a type
decision: the solid and striped rules use the sum of these two categories.
Searching it as a type parameter would duplicate equivalent predictions.
The separate black-ball thresholds in number recognition were searched in
stage two.

Configurations were ranked first by correct full-number predictions, then by
correct type predictions, then by correct oracle-number predictions. All
counts used Table1--Table6 only. Remaining ties favoured the smaller departure
from the original parameters. The departure summed absolute threshold changes
normalised by their search step sizes, plus $K-2$ and one unit for enabling
cloth filtering. The hue boundary at 165 used a five-unit normalisation even
though its upper alternative was 168; the body-area threshold used absolute
percentage-point change. Final ties favoured the earlier configuration ID.
This tie rule prefers less change when measured performance is equal; it
does not establish that the chosen thresholds are uniquely optimal. A
separate best-type ranking reversed the first two accuracy criteria.

% Generated by experiments/writeRecognitionReport.py from saved CSVs.
\begin{table}[htbp]
\centering
\small
\caption{Effect of K and cloth filtering on Table1--6. Default columns use the original thresholds. Tuned columns select the highest-ranked stage-one configuration within each K/filter pair; colour thresholds remain fixed.}
\label{tab:recognition-search-k}
\begin{tabular}{rcrrrrr}
\toprule
K & Cloth & Default type & Default no. & Tuned type & Tuned no. & ID \\
\midrule
2 & Off & 76/86 & 67/86 & 80/86 & 71/86 & 70 \\
2 & On & 77/86 & 69/86 & 79/86 & 71/86 & 249 \\
3 & Off & 73/86 & 61/86 & 81/86 & 69/86 & 370 \\
3 & On & 77/86 & 67/86 & 80/86 & 70/86 & 610 \\
4 & Off & 73/86 & 65/86 & 79/86 & 71/86 & 820 \\
4 & On & 76/86 & 67/86 & 79/86 & 70/86 & 1060 \\
5 & Off & 73/86 & 65/86 & 79/86 & 71/86 & 1179 \\
5 & On & 76/86 & 68/86 & 79/86 & 71/86 & 1359 \\
6 & Off & 74/86 & 66/86 & 81/86 & 73/86 & 1510 \\
6 & On & 76/86 & 67/86 & 81/86 & 72/86 & 1750 \\
\bottomrule
\end{tabular}
\end{table}
\FloatBarrier


The default-threshold columns show that increasing K alone did not improve
number recognition. Without cloth filtering, the original K=2 configuration
correctly numbered 67 balls, compared with 61, 65, 65 and 66 for K=3--6.
However, after joint type-threshold tuning, K=6 without filtering achieved
73 correct numbers and 81 correct types. Thus the effect of K depended on
how the resulting clusters were interpreted; the experiment did not support
simply increasing K while leaving all thresholds fixed.

\subsubsection{Stage two: grouped colour-threshold search}

To limit the much larger colour search, a shortlist was formed from the
top three distinct K/filter pairs under the full-number ranking and the
top three distinct K/filter pairs under the type-first ranking. The original
K=2 baseline, with and without cloth filtering, was also included. Exact
duplicate seed configurations were removed. This run produced six seeds.
Distinctness within each ranking avoided filling the shortlist with several
nearly identical threshold settings for the same K/filter pair.

Each seed retained its K, cloth setting and stage-one type thresholds.
Colour thresholds were searched in three consecutive groups:

\begin{enumerate}
\item Warm colours: red's upper hue boundary $\{4,6,8\}$;
orange's upper boundary $\{12,15,18\}$; yellow's upper boundary
$\{30,35,40\}$; red's wrap-around lower boundary
$\{168,170,172\}$; and the red/maroon brightness split
$\{130,150,170\}$. The nominal Cartesian product contained 243 combinations.
\item Green, blue and purple: green's lower boundary
$\{40,45,50\}$; green's upper boundary $\{85,90,95\}$;
blue's lower boundary $\{95,100,105\}$; the blue/purple boundary
$\{125,130,135\}$; and purple's upper boundary
$\{160,165,168\}$. This also gave 243 nominal combinations.
\item Neutral colours and body selection: black's saturation split
$\{40,60,80\}$; black's maximum brightness $\{70,90,110\}$;
and minimum eligible body-cluster area $\{1,2,5\}\%$.
This gave 27 nominal combinations.
\end{enumerate}

The intervals retained the original logic: red occupies both ends of the
hue scale; orange and yellow follow its low-hue interval; green and blue
have separate lower bounds; blue and purple share a boundary. Hue intervals
must remain ordered, with strictly positive gaps between yellow and green,
green and blue, and purple and high-hue red. Combinations violating these
constraints were rejected. Existing unknown-colour gaps were therefore
preserved rather than forcing every hue into a ball colour.

For each group, the current configuration was retained as a candidate. The
group's best configuration under the same full-number ranking became the
starting point for the next group. Two complete passes revisited warm,
cool and neutral thresholds after later groups had changed. This was a
sequence of exhaustive searches within small groups, not an exhaustive
Cartesian search over every colour parameter jointly. It can miss a useful
combination requiring simultaneous changes across groups, and it does not
prove global optimality. Stage-one type thresholds were not revisited after
colour tuning. Their remaining interactions with colour thresholds are
another limitation of the staged search.

The actual second stage evaluated 5,526 valid combinations. Together with
stage one this gave 7,326 scored trials. Trials repeated across groups,
seeds or passes were retained in the CSV audit trail; this count is not
the number of unique parameter vectors. Each trial recorded its parent
configuration and stage so the selection sequence could be reconstructed.
The minimum eligible cluster percentage was included for completeness, but
with at most six clusters the largest non-light cluster often already
exceeds all tested minimum areas. A parameter with no effect on these
images should not be interpreted as precisely estimated.

\subsubsection{Selected configuration and measured improvement}

% Generated by experiments/writeRecognitionReport.py from saved CSVs.
\begin{table}[htbp]
\centering
\small
\caption{Baseline and selected recognition parameters. Other colour parameters retain their original values.}
\label{tab:recognition-search-parameters}
\begin{tabular}{lrr}
\toprule
Parameter & Baseline & Selected \\
\midrule
K & 2 & 6 \\
Cloth filter (0=off) & 0 & 0 \\
Light minimum V & 170 & 170 \\
Light maximum S & 165 & 125 \\
Cue minimum light (\%) & 85 & 85 \\
Solid minimum non-light (\%) & 70 & 80 \\
Blue/purple hue boundary & 130 & 125 \\
\bottomrule
\end{tabular}
\end{table}
\FloatBarrier

% Generated by experiments/writeRecognitionReport.py from saved CSVs.
\begin{table}[htbp]
\centering
\small
\caption{Exploratory recognition scores. Tuning uses Table1--6; the final column is full-number recognition on AllBalls, reported separately.}
\label{tab:recognition-search-summary}
\begin{tabular}{lrrrrrr}
\toprule
Configuration & Type & Type acc. & Number & Number acc. & Oracle & Ref. no. \\
\midrule
Original baseline & 76/86 & 88.37\% & 67/86 & 77.91\% & 70/86 & 16/16 \\
Baseline + cloth filter & 77/86 & 89.53\% & 69/86 & 80.23\% & 72/86 & 15/16 \\
Best K/type grid & 81/86 & 94.19\% & 73/86 & 84.88\% & 74/86 & 15/16 \\
Best after colour search & 81/86 & 94.19\% & 74/86 & 86.05\% & 75/86 & 16/16 \\
\bottomrule
\end{tabular}
\end{table}
\FloatBarrier


The highest-ranked configuration used K=6, no cloth filtering,
$V\geq170$ and $S\leq125$ for light clusters, an 85\% light threshold
for the cue ball, and an 80\% non-light threshold for solids. Relative to
the original rules, its blue/purple hue boundary moved from 130 to 125.
Other colour thresholds remained unchanged. It was configuration 2153
in the saved run. The same configuration won both the full-number and
type-first rankings after tie-breaking, so these names do not refer to
two distinct selected models.

On the 86 matched tuning balls, type correctness increased from 76 to 81:
$88.37\%$ to $94.19\%$, a gain of 5.81 percentage points. Full-number
correctness increased from 67 to 74: $77.91\%$ to $86.05\%$, a gain of
8.14 percentage points. Five type errors and seven number errors were
corrected relative to the baseline, with no previously correct predictions
becoming incorrect in either category on these matched candidates.

The best stage-one configuration already produced 73 correct numbers.
Colour tuning added one further correct number on Table1--Table6 and
restored the stage-one reference score on AllBalls from 15/16 to 16/16.
The latter restoration was observed after selection, not used to choose
the hue boundary. Oracle-number correctness increased from 70/86 for the
baseline to 75/86 for the selected configuration. It remained only one
ball above full-number correctness, indicating that supplying the correct
type alone resolved little of the remaining number-recognition error.

The cloth filter improved the original thresholds slightly, but was not
part of the selected overall configuration. Its benefit was therefore
conditional on K and the thresholds. These results support keeping cloth
filtering optional rather than automatically adding it to the main pipeline.

% Generated by experiments/writeRecognitionReport.py from saved CSVs.
\begin{table}[htbp]
\centering
\small
\caption{Per-image counts before and after the recognition search. All comparisons use the same matched candidates.}
\label{tab:recognition-search-images}
\begin{tabular}{lrrrrrr}
\toprule
Image & GT & Matched & Base type & Best type & Base no. & Best no. \\
\midrule
Table1.jpeg & 16 & 14 & 8 & 11 & 6 & 9 \\
Table2.jpeg & 16 & 16 & 15 & 16 & 13 & 15 \\
Table3.jpeg & 16 & 15 & 14 & 14 & 14 & 14 \\
Table4.jpeg & 15 & 14 & 14 & 14 & 13 & 13 \\
Table5.jpeg & 14 & 14 & 14 & 14 & 11 & 12 \\
Table6.jpeg & 13 & 13 & 11 & 12 & 10 & 11 \\
\bottomrule
\end{tabular}
\end{table}
\FloatBarrier


Improvements were unevenly distributed. On Table1, correct type predictions
increased from 8/14 to 11/14 and correct numbers from 6/14 to 9/14. On
Table2, correct numbers increased from 13/16 to 15/16. Tables5 and 6 each
gained one correct number, while Tables3 and 4 retained their baseline
number counts. This per-image view prevents the pooled gain from being
mistaken for an equally large improvement on every scene.

Using all 90 annotations as the denominator, correct detected-and-numbered
balls increased from 67/90 ($74.44\%$) to 74/90 ($82.22\%$).
Detection itself was unchanged; the four missed balls could not be recovered
by recognition tuning. The 16/16 reference result also matched, rather than
exceeded, the original baseline on AllBalls.

% Generated by experiments/writeRecognitionReport.py from saved CSVs.
\begin{table}[htbp]
\centering
\small
\caption{Selected configuration's type confusion matrix on Table1--6. Rows are ground truth and columns are predictions.}
\label{tab:recognition-search-confusion}
\begin{tabular}{lrrrr}
\toprule
Actual / predicted & White & Solid & Striped & Unknown \\
\midrule
White & 6 & 0 & 0 & 0 \\
Solid & 0 & 42 & 1 & 0 \\
Striped & 3 & 1 & 33 & 0 \\
\bottomrule
\end{tabular}
\end{table}
\FloatBarrier


The selected configuration still made five type errors and twelve full-number
errors among matched tuning detections. For example, some striped balls were
predicted as white because their visible light region dominated the crop.
Some purple balls were assigned another colour or returned unknown; two black
balls were assigned yellow. These examples show that tuning global thresholds
does not resolve all effects of ball orientation, shadows, highlights and
the selection of a single representative body cluster. They also explain
why most remaining number errors persisted when the true type was supplied.

\begin{figure}[htbp]
\centering
\includegraphics[width=\linewidth]{recognition_search_example.png}
\caption{One matched candidate from Table1, compared under the original baseline,
the cloth-filter baseline, the best stage-one configuration and the final
selected configuration. Only final masked ball crops are shown. Each panel
reports K, cloth setting, predicted type and number. Green text denotes a
correct prediction and red denotes an incorrect prediction against the
matched annotation. Identical winning configurations are displayed once.}
\label{fig:recognition-search-example}
\end{figure}
\FloatBarrier

\subsubsection{Reproducibility, outputs and remaining limitations}

The search is run from the repository root with
\code{./build-local/BallExperiments --recognition-grid} after building the
\code{BallExperiments} CMake target. Search ranges and original defaults are
defined in \code{experiments/recognitionSearchConfig.h}. An optional output
directory argument allows separate runs to be preserved.

Outputs are written under \code{output/experiments/recognition_grid_search/}.
The type and colour leaderboards contain every evaluated configuration,
its complete parameter vector, ranking metrics, stage and parent identifier.
The K-comparison table records both the original thresholds and the best
stage-one thresholds for each K/filter pair. Other CSV files contain
the seed shortlist, finalist configurations, pooled and per-image metrics,
candidate-level predictions, type and number confusion matrices, detection
counts, mask sizes and cached cluster features. Unmatched candidates have
blank correctness fields rather than being silently counted as correct or
incorrect classifications. A separate image is saved for each candidate,
with the baseline and distinct finalist configurations placed side by side.

Report tables and the diagnostic figure are local snapshots produced by
\code{python3 experiments/writeRecognitionReport.py}; the report can compile
without the experiment outputs. Rerunning the search does not automatically
update those snapshots. The surrounding interpretation must also be reviewed
if parameters or measured results change.

The main limitation is selection on a small, correlated image set. Thousands
of threshold trials can fit properties specific to these balls, their
orientations and the shared illumination. The measured gains should therefore
motivate evaluation on fresh photographs before replacing the baseline.
Neither the pooled tuning result nor the AllBalls reference check establishes
performance on independent data. The search was also bounded to five K values,
small threshold grids and two grouped colour passes; it establishes the best
observed configuration under that procedure, not the best possible recognition
system. No new production settings were applied as part of this experiment.
